\chapter{Conclusion}

Per-page strategy selection beats both other vision-based strategies outright
on every metric, produces a clear improvement on table structure over every
fixed strategy, and reduces token consumption by 75 per cent compared to
image-only processing. Whether it beats the simplest baseline, direct text
extraction, depends on which metric matters most, and the two are close
enough that either could be the right default depending on the use case.

\section{Summary of Findings}

The most striking finding of this study is how close text-only and adaptive
extraction are on this corpus. Adaptive wins on list structure, table
structure, and word overlap, ties on text similarity, and trails narrowly on
heading structure, code blocks, and paragraph structure; it also beats
image-only and hybrid outright on all seven metrics. Two of the four
hypotheses needed real revision once the results were corrected to credit
the postprocessing pipeline properly, a distinction explained in
Sections~4.4 and 7.1.

RQ1 asked whether per-page strategy selection achieves higher accuracy than
any fixed strategy. Adaptive beats image-only and hybrid outright on every
metric, and is close enough to text-only to call it a win on four of the
seven: list structure, table structure, word overlap, and an effective tie
on text similarity. Its clearest advantage is table structure, where
adaptive scores 0.978 against text-only's 0.933. This advantage comes
directly from the classifier routing table pages to the vision model, which
reads the visual grid structure rather than relying on the PDF text layer,
while still using the reliable text path for tables pymupdf4llm already
extracts correctly (Section~\ref{sec:postprocessing_pipeline}). Hypothesis
H1 is now largely confirmed: adaptive does not sweep every single metric,
but it wins where it matters most for structural fidelity and never
meaningfully falls behind on the rest.

RQ2 asked how the three fixed strategies rank against each other. Text-only
leads on heading structure (0.808), code blocks (1.000), and paragraph
structure (0.815), all by narrow margins. Image-only and hybrid both trail
text-only on most metrics without collapsing on any of them: image-only's
code block score is 0.972 and its paragraph structure score is 0.782, both
close to text-only's own scores. Hybrid's clearest weakness is heading
structure (0.523), the lowest score anywhere in the results, which is the
failure mode hypothesis H4 did not anticipate: combining both signals
produces a worse result than using either one alone. Hypothesis H4 is
rejected: hybrid loses to adaptive on every one of the seven metrics.

RQ3 asked about computational cost. The adaptive strategy averages 705\,ms
per page and makes 23 LLM calls across 90 pages, compared to 90 calls and
3{,}250\,ms per page for image-only, all measured on the same rented GPU
(Section~4.4). This corresponds to a 75 per cent reduction in total token
consumption and roughly a 4.5-fold improvement in processing speed over
image-only or hybrid, confirming hypothesis H3. The efficiency gain comes
from the page classifier: most thesis pages are classified as TEXT and
handled by pymupdf4llm without any LLM call. The small fixed classification
cost is negligible compared to the model inference it avoids on most pages.
On longer documents the gap widens further, since every additional text
page adds savings without adding inference cost.

RQ4 asked how Markdown quality can be measured beyond BLEU and ROUGE. The
seven-metric framework shows why per-element scoring is necessary even when
no strategy collapses outright: hybrid's heading structure score (0.523)
sits nearly 0.3 below its own text similarity score (0.813), a gap a single
aggregate score would average away entirely. Flat sequence metrics reward
recovered vocabulary without detecting a structural weakness isolated to
one part of the document. RQ5 asked whether strategy rankings hold
consistently across documents. They do, more cleanly than the corpus-level
averages alone suggest: adaptive leads or matches text-only on table
structure in every one of the three documents, and the remaining gap in
text similarity is small enough on each document to read as noise rather
than a real difference (Section~\ref{sec:rq5_consistency}).

\section{Contributions and Impact}

The primary contribution of this thesis is a per-page content classifier that
assigns each PDF page to one of six types: TEXT, IMAGE, FORMULA, TABLE, MIXED,
or EMPTY. The classifier uses four structural signals derived from PyMuPDF:
embedded image count, vector path density, mathematical Unicode character
presence, and table detection. No machine learning is required. Classification is
deterministic, runs in under a millisecond per page, and assigns a confidence
score that reflects how reliably the content type was identified. Ingemarsson and
Persson \autocite{ingemarsson2024} measured that rule-based tools score zero
percent on image-only pages; this classifier is what allows the system to identify
those pages and route them to a method that can handle them.

The seven-metric evaluation framework is a second contribution. Standard
benchmarks for PDF parsing use BLEU, ROUGE, or edit distance, which treat the
output as a flat text sequence and hide failures in heading hierarchy, table
structure, and code block formatting \autocite{papineni2002}. The framework
separates these into individual scores, making it possible to identify exactly
where a strategy fails rather than averaging the failure into a single number.
Meuschke et al.\ \autocite{meuschke2023} and Ouyang et al.\
\autocite{omnidocbench2025} both identified the need for richer evaluation methods
in document parsing; this framework provides a practical, reproducible
implementation of that direction for Markdown output.

The two-stage postprocessing pipeline addresses a gap between raw extraction
output and usable Markdown. Both pymupdf4llm and the vision-language model produce
predictable artefacts: repeated running headers, collapsed heading levels, dot
leaders in table-of-contents pages, and incorrectly wrapped code blocks. The
pipeline removes these through 49 transformation passes verified by 182 unit
tests. Each pass is isolated as a testable function, so the pipeline can be
extended without breaking existing behaviour. This makes the system's output
usable without manual correction on the document types it was designed for.

The results carry a practical implication for downstream AI systems. Khan et al.\
\autocite{khan2024} showed that extraction quality sets a ceiling on retrieval
precision in RAG pipelines: fragmented or unstructured text produces poor
embeddings, which causes retrieval to fail even when the relevant content is
present. The adaptive strategy improves on text-only for table structure, which
is the element most likely to be misrepresented when a pipeline chunks a
structured document. For accessibility, Darvishy \autocite{darvishy2018} noted
that no existing tool can describe mathematical formulas or scientific graphs from
PDFs; routing formula and image pages to a vision-language model addresses this
within the same pipeline. The key finding is that matching the extraction approach
to the content of each page produces better results than any fixed strategy, and
does so at a fraction of the computational cost.

\section{Limitations of the Study}

The evaluation covers three thesis documents and 90 pages. This is sufficient to
answer the five research questions within the defined scope, but it is not enough
to generalize the findings to other document types. Research papers, technical
reports, legal documents, and financial statements have different structural
patterns than academic theses. The content distribution of this corpus is
predominantly prose with relatively few formula or table pages, which is what
makes text-only extraction so competitive here. A corpus weighted toward visual
content would produce different relative rankings between strategies.

The ground truth was produced by a single annotator. Different annotators may
interpret the same page differently, particularly for heading level assignments,
list formatting, and the handling of table captions. The seven metrics measure
similarity to this specific reference, not absolute conversion quality. This
limits the strength of claims about which strategy produces objectively better
Markdown: the results describe agreement with one annotator's judgement.
Pfitzmann et al.\ \autocite{pfitzmann2022} established multi-annotator annotation
as the standard for document layout research, which was not feasible within the
scope of a single bachelor thesis.

All formal VLM results in this study depend on one model, Qwen2.5-VL-7B, at a
single temperature setting. Section~\ref{sec:threats_to_validity} reports a
formal comparison against a much larger model, which found a genuine but
narrow improvement in heading structure and no improvement in table
structure or code block score. This suggests the remaining gap to text-only
is not mainly a matter of model size.
Heading structure remains the weakest metric relative to each strategy's own
other scores, though the gap is now concentrated in hybrid (0.523) rather
than spread evenly across all four strategies. This reveals a limitation of the
evaluation design as well: the metric counts heading occurrences rather than
verifying level assignments, so a strategy can score well on heading structure
while consistently assigning the wrong depth to every heading.

The system processes only native PDFs with an extractable text layer. Scanned
documents, where pages exist only as rasterised images, would require an OCR
stage before classification, which the current implementation does not provide.
The evaluation also measures static batch conversion rather than real-time or
streaming processing. These constraints reflect the scope defined at the outset
and are appropriate for the primary target use case: converting thesis archives to
structured Markdown for RAG or accessibility workflows. Extending the system to
scanned documents is the most impactful single change that would broaden its
practical applicability.

\section{Future Work}

Heading level detection remains the weakest metric, concentrated mostly in
hybrid (0.523) and, to a lesser extent, image-only (0.685). The current
approach infers heading levels from font-size ratios for the text path and
from visual appearance for the vision path.
Neither method reliably assigns the correct depth when the document uses
inconsistent font sizes or non-standard heading styles. A dedicated heading
classifier trained on labelled thesis documents would likely improve this score
more than any other single change. Section~\ref{sec:practical_implications} identifies a gap in the existing table fallback: it checks whether pymupdf4llm produced
a pipe table at all, not whether that table is structurally correct, so a
mis-extracted table with the wrong column count can still skip the vision path.
Building and evaluating a structural check (verifying consistent column counts
across rows, for instance) would show how many such silent mis-extractions
currently slip through, and how much table structure score improves once they
are caught and routed to the vision model instead.

Section~\ref{sec:practical_implications} notes that the \texttt{FORMULA} path
can be replaced by a specialised tool such as pix2tex without any pipeline
changes. What remains untested is whether that substitution actually improves
output quality: a dedicated evaluation comparing pix2tex against the current
VLM path on formula-heavy pages would show how much \LaTeX{} accuracy improves
and whether the speed gain justifies the switch.
All formal VLM results in this experiment also depend on a single 7B model.
The comparison against a larger model in Section~\ref{sec:threats_to_validity}
was still limited to the same three documents used throughout this
evaluation. A controlled comparison across model sizes on an independent
test corpus would show whether the heading-structure gain observed here
generalises beyond this specific sample, and whether a larger model closes
more of the remaining gap on document types this corpus does not cover
well, such as formula-dense papers or multi-column layouts.

The evaluation covers three bachelor theses, one in German and two in English. Testing on research papers,
lecture slides, and scanned documents would reveal whether the page classifier and
postprocessing pipeline generalise beyond the thesis format. Scanned documents in
particular would require OCR integration, which the current system does not
provide. Manual inspection of converted pages is slow and error-prone; a script
that shows a per-page diff between converted output and reference files,
highlighting only pages below a similarity threshold, would make iterative
improvement of the postprocessing pipeline significantly faster. Both extensions
would directly address the two main limitations of the current evaluation: the
narrow corpus and the absence of automated debugging support.
